\notechapter{N-20221227}{Category Theory Continued and the Question Topology}

\section{More examples of categories}

This note continues the category-theory examples.  Abelian groups with group homomorphisms form \(\mathbf{Ab}\).  Modules over a ring \(A\) form \(\mathbf{Mod}_A\).  Rings with unital ring homomorphisms form \(\mathbf{Rings}\).  Topological spaces with continuous maps form \(\mathbf{Top}\), whose isomorphisms are homeomorphisms.

A partially ordered set \((S,\ge)\) can be regarded as a category: the objects are elements of \(S\), and there is exactly one morphism \(x\to y\) iff \(x\ge y\).  Reflexivity gives identities, transitivity gives composition, and antisymmetry ensures isomorphisms are trivial.

\section{Subcategories and functors}

A subcategory \(\mathcal A\subseteq\mathcal C\) has some objects and some morphisms of \(\mathcal C\), closed under identities and composition.  A functor
\[
  F:\mathcal C\to\mathcal D
\]
sends objects to objects and morphisms to morphisms, preserving identities and composition:
\[
  F(1_A)=1_{F(A)},\qquad F(g\circ f)=F(g)\circ F(f).
\]
A standard functoriality square says that the image of a composable path under \(F\) is still composable and that the composite maps to the composite.

\section{Forgetful functors}

A forgetful functor forgets structure.  For example,
\[
  \mathbf{Grp}\to\mathbf{Set}
\]
sends a group to its underlying set and a homomorphism to its underlying function.  Similarly,
\[
  \mathbf{Top}\to\mathbf{Set}
\]
forgets open sets, and \(\mathbf{Vect}_k\to\mathbf{Set}\) forgets vector-space operations.

A forgetful functor is not innocent: once structure is forgotten, many questions can no longer be answered.  A bijection of underlying sets need not be an isomorphism of groups, vector spaces, or topological spaces.

\section{A topology generated by a distance-like question}

The second half moves into topology.  The page considers a collection of subsets and asks what topology it generates.  Given a basis-like family \(\mathcal B\), the topology generated by it consists of arbitrary unions of finite intersections of elements of \(\mathcal B\).

The displayed rational-looking function and its graph on the page are used as a reminder: topology is not only about familiar Euclidean open intervals.  One can impose a topology by declaring certain sets to be basic neighborhoods.


\section{Posets as categories, with an example}

A partially ordered set becomes a category in a very economical way.  There is a unique morphism \(x\to y\) exactly when \(x\le y\).  Identity morphisms come from reflexivity \(x\le x\), and composition comes from transitivity: if \(x\le y\) and \(y\le z\), then \(x\le z\).

This example is important because it shows that a category need not look like a collection of algebraic objects and homomorphisms.  A category can encode logic or order.  Meets and joins in a poset become limits and colimits in the corresponding category.

\section{Functorial thinking}

A functor is not merely a map between categories.  It is a structure-preserving translation.  If \(F:\mathcal C\to\mathcal D\), then for every composable pair
\[
  A\xrightarrow{f}B\xrightarrow{g}C
\]
we require
\[
  F(g\circ f)=F(g)\circ F(f).
\]
This condition says that doing a construction after composing maps gives the same result as first applying the construction to each map and then composing.

Examples clarify the point.  The fundamental group construction is a functor from pointed topological spaces to groups.  Homology is a functor from spaces to graded abelian groups.  The forgetful functor from groups to sets forgets multiplication but still sends homomorphisms to functions.

\section{Generated topology}

The topology part of the note concerns how open sets can be generated from a chosen collection.  If \(\mathcal B\) is intended as a basis, two conditions are expected: every point lies in some basis element, and intersections of basis elements are locally refined by basis elements.  Then every open set is a union of basis elements.

If one starts with a subbasis \(\mathcal S\), the generated topology consists of arbitrary unions of finite intersections of elements of \(\mathcal S\).  This explains the phrase “generated by”: we first force the chosen sets to be open, then close under the axioms of topology.

\section{Why category theory appears beside topology}

Topological constructions are full of universal properties.  Product topology is the coarsest topology making the projection maps continuous.  Quotient topology is the finest topology making the quotient map continuous.  These statements are categorical: they characterize objects by how maps into or out of them behave.

The passage from categories to topology is natural.  Topology studies spaces through continuous maps, not only through open sets, and this is precisely categorical thinking.

\section{The larger pattern}

Category theory and topology meet through structure-preserving maps.  A functor preserves categorical structure; a continuous map preserves open-set structure.  A forgetful functor deliberately discards structure.  This is why the note naturally moves from categories to topology: both ask what data must be preserved for two objects to be considered the same.

\section{Posets as categories and logic}

A poset viewed as a category has at most one morphism between any two objects.  In this setting, categorical statements become order-theoretic statements.  A product is a meet, a coproduct is a join, an initial object is a least element, and a terminal object is a greatest element.

For example, in the poset of open subsets of a space ordered by inclusion, intersections behave like meets and unions behave like joins.  This is why topology and order theory are so close: open sets form a lattice, and categorical language packages its universal properties.

\section{Forgetful functors and free constructions}

A forgetful functor loses structure, but it often has a left adjoint that freely adds structure.  For instance,
\[
  U:\mathbf{Grp}\to\mathbf{Set}
\]
forgets the group operation.  Its left adjoint sends a set \(S\) to the free group \(F(S)\).  The adjunction says that a set map
\[
  S\to U(G)
\]
corresponds uniquely to a group homomorphism
\[
  F(S)\to G.
\]
This is the categorical form of "define a homomorphism out of a free object by choosing images of generators."

\section{Initial topology generated by questions}

The topology generated by a family of maps can be understood as an initial topology.  Suppose we have functions
\[
  f_i:X\to Y_i
\]
where each \(Y_i\) is already topological.  The initial topology on \(X\) is the coarsest topology making all \(f_i\) continuous.  Its subbasis is
\[
  f_i^{-1}(U),\qquad U\subseteq Y_i\text{ open}.
\]
This realizes the idea of a topology generated by observations: each map \(f_i\) observes \(X\), and the generated topology is exactly the open-set structure forced by those observations.

\section{Product topology as an initial topology}

The product topology is the initial topology for the projection maps
\[
  \pi_X:X\times Y\to X,\qquad \pi_Y:X\times Y\to Y.
\]
A map \(h:Z\to X\times Y\) is continuous iff the coordinate maps \(\pi_X\circ h\) and \(\pi_Y\circ h\) are continuous.  This is a categorical universal property, not a visual fact about rectangles.

The rectangle basis \(U\times V\) appears because
\[
  U\times V=\pi_X^{-1}(U)\cap\pi_Y^{-1}(V).
\]
Thus the usual construction of product open sets is forced by the initial-topology viewpoint.

\section{Quotient topology as a final topology}

Dually, quotient topology is a final topology.  For a quotient map
\[
  q:X\to Y,
\]
the topology on \(Y\) is the finest topology making \(q\) continuous, and a map \(f:Y\to Z\) is continuous iff \(f\circ q:X\to Z\) is continuous.

Thus product and quotient are opposite-looking constructions:
\[
\begin{array}{c|c|c}
\text{initial topology} & \text{maps into known spaces} & \text{coarsest topology}\\
\text{final topology} & \text{maps from known spaces} & \text{finest topology}
\end{array}
\]
This is the precise categorical version of the topology examples in the note.

\section{Functors from topology to algebra}

The fundamental group is functorial.  More generally, algebraic topology studies functors such as
\[
  \pi_1:\mathbf{Top}_*\to\mathbf{Grp},\qquad
  H_n:\mathbf{Top}\to\mathbf{Ab}.
\]
A continuous map of spaces induces a homomorphism of groups or abelian groups.  This turns geometric questions into algebraic ones.

The power of this move is that algebraic invariants can obstruct homeomorphisms.  If two spaces have non-isomorphic fundamental groups or homology groups, they cannot be the same topological space.

\section{Initial and final structures in topology}

The category-theory material is not separate from topology.  Posets show that categories can encode order and logic.  Functors express structure-preserving translation.  Generated topologies, products, and quotients are governed by universal properties.  Algebraic topology then uses functors to extract algebra from spaces.  The common theme is: understand objects through the maps they admit.
